\section{结论}
本文根据题给环境、经验物性与半径数据，建立并整理了固定圆柱和收缩圆柱两类有效热质传递模型。通过守恒型有限体积离散、自适应BDF积分和连续阈值定位，获得四问局部分布与干燥终点。主要结论如下。

\begin{enumerate}
\item 第一问1800 s时，中心、表面温度分别为33.5753、36.7856 $^\circ$C，含水率分别为2.5500、1.5102 kg/kg。热量逐步向内传播，而明显失水首先集中于表层，预热完成不等价于干燥完成。
\item 第二问采用统一变物性后，3 h时中心和表面温度已达49.8495、49.9664 $^\circ$C，但含水率仍为1.7662、1.0081 kg/kg。温度接近稳定与全域水分达标具有不同时间尺度，需以含水率阈值独立判断结束。
\item 第三问在固定半径、后期环境取3--4 h均值的条件下，最大含水率降至0.15 kg/kg的临界时长为57.4740 h，最后达标位置为中心。表面在较早阶段已经很干，不能作为全部材料的停止依据。
\item 第四问考虑均匀径向收缩并使用新物性，临界时长为51.0920 h、半径为1.2000 cm。与第三问相比减少约6.3820 h，即11.10\%；这一差异包含物性和几何的共同影响。
\item 独立级数、网格加密和容差收紧支持现有数值结果的稳定性；后期环境选择造成约0.3 h的差异，提示实际应用应优先补充后期环境与材料边界数据，而非仅增加输出小数位。
\end{enumerate}

上述结果已与既有表格和计算摘要逐项对照。本文使用的主算法保持不变，未为改善论文图像或结论而重新计算或调整数值；所有比较均来源于此前保存的主计算及验证记录。

\begin{table}[H]\centering
\caption{四问核心结果汇总}\label{tab:final}
\begin{tabularx}{\linewidth}{lX}
\toprule 项目 & 核心结果\\\midrule
问题一，1800 s & $T_c/T_s=33.5753/36.7856\,^\circ$C；$C_c/C_s=2.5500/1.5102$\\
问题二，3 h & $T_c/T_s=49.8495/49.9664\,^\circ$C；$C_c/C_s=1.7662/1.0081$\\
问题三，临界 & 57.4740 h；固定半径2.0000 cm\\
问题四，临界 & 51.0920 h；当前半径1.2000 cm\\
\bottomrule
\end{tabularx}
\end{table}

\subsection*{附录A\quad AI使用说明}
\addcontentsline{toc}{section}{附录A AI使用说明}
本文使用AI辅助整理既有代码、计算结果、模型推导、文献元数据及LaTeX排版；数值来源为已有Python求解与验证记录，本轮未重新运行模型。AI还辅助制作流程示意及图表改绘建议。模型假设、原始数据与引用依据均在文中或项目记录中明示；正式提交前由参赛成员独立复核并按赛事当届规定完善使用声明。
